\documentclass[10pt,a4paper]{article}
\usepackage{fontspec,amsmath,amssymb,graphicx}
\usepackage[margin=18mm]{geometry}
\setmainfont{Times New Roman}
\setlength{\parindent}{0pt}
\setlength{\parskip}{5pt}
\allowdisplaybreaks
\sloppy
\emergencystretch=2em
\begin{document}
\section*{BT\_4\_5\_3\_01 — Lời giải bổ sung}
\subsection*{Kiến thức và trình tự giải}
\textbf{Bước 1 — Loại đạo hàm cấp một bằng tính đạo hàm đầy đủ.} Với $u_t=Ku_{xx}+Au_x+Bu$, đặt $u=e^{-hx}v$, $h=A/(2K)$. Khi đó


\[
u_t=e^{-hx}v_t,\quad u_x=e^{-hx}(v_x-hv),\quad
u_{xx}=e^{-hx}(v_{xx}-2hv_x+h^2v).
\]


Thay vào PDE, hệ số $v_x$ là $A-2Kh=0$, hệ số $v$ là $Kh^2-Ah+B=B-A^2/(4K)$. Vì thế $v_t=Kv_{xx}+Cv$, $C=B-A^2/(4K)$; có thể tiếp tục đặt $v=e^{Ct}z$ để $z_t=Kz_{xx}$.

\textbf{Bước 2 — Đổi dữ liệu và biên cùng phép đổi ẩn.} $v(x,0)=e^{hx}f(x)$. Dirichlet thuần nhất giữ nguyên; Neumann biến thành $v_x-hv=0$, nên không giữ nguyên một cơ sở cosine nếu $A\ne0$. Với Dirichlet:


\[
u=e^{-hx}\sum_{n\ge1}\left(\frac2L\int_0^Le^{hy}f(y)\sin\frac{n\pi y}Ldy\right)
e^{[C-K(n\pi/L)^2]t}\sin\frac{n\pi x}L.
\]


\textbf{Bước 3 — Tính hệ số mũ–đa thức.} Dùng nguyên hàm


\[
\int e^{hy}\sin ky\,dy
=\frac{e^{hy}(h\sin ky-k\cos ky)}{h^2+k^2}.
\]


Với $y e^{hy}$ hoặc $y(1-y)e^{hy}$, lấy đạo hàm theo $h$ của tích phân đã tính trên cận cố định; trình bày đạo hàm để tìm hệ số thực. Cách này tránh để lại một tích phân sơ cấp chưa thực hiện.

\textbf{Bước 4 — Hệ số không gian biến thiên.} Với $x^2u_{xx}$ trên $x>0$, đặt $z=\ln x$: nếu $U(z,t)=u(e^z,t)$ thì $u_x=U_z/x$, $x^2u_{xx}=U_{zz}-U_z$. Đổi miền và biên sang biến $z$ rồi dùng bước 1. Biên không thuần nhất nên tìm trạng thái dừng trước khi khai triển.

\textbf{Bước 5 — Duy nhất và ổn định.} Với hiệu nghiệm Dirichlet, $I'=-K\int w_x^2+2BI\le2BI$ suy $w=0$ khi $I(0)=0$. Giới hạn phụ thuộc các số $C-K\lambda_n$ có hệ số khác không: âm thì giảm, bằng không thì tồn tại mode dừng, dương thì tăng. Phải kiểm tra tính thống nhất giữa miền PDE và vị trí biên trước khi dùng công thức.
\subsection*{Lời giải áp dụng cho bài}
\textbf{Bước 1. Phép đổi ẩn cho hai biên Dirichlet.}
$u=e^{-Ax/(2K)}v$ cho $v_t=Kv_{xx}+(B-A^2/(4K))v$, dữ liệu $e^{Ax/(2K)}f$, biên $0$. Cơ sở là $\sin nx$ vì $L=\pi$.
\textbf{Bước 2. Câu (i), $A=-1,B=0,K=1,f=1$.}
Đặt $h=-1/2$, dùng nguyên hàm $e^{hx}(h\sin nx-n\cos nx)/(h^2+n^2)$:
\[c_n=\frac2\pi\int_0^\pi e^{-y/2}\sin ny\,dy
=\frac{2n(1-(-1)^ne^{-\pi/2})}{\pi(n^2+1/4)}.\]
Suy
\[u=e^{x/2}\sum_{n\ge1}c_ne^{-(n^2+1/4)t}\sin nx.\]
\textbf{Bước 3. Câu (ii), $A=0,B=-1,K=1,f=x$.}
Tích phân từng phần:
\[c_n=\frac2\pi[-y\cos ny/n+\sin ny/n^2]_0^\pi
=\frac{2(-1)^{n+1}}n,\qquad
u=\sum_{n\ge1}\frac{2(-1)^{n+1}}n e^{-(n^2+1)t}\sin nx.\]
\textbf{Bước 4. Chứng minh duy nhất bằng năng lượng.}
Hiệu $w$ có dữ liệu $0$ và biên $0$. Tính
\[I'=K[ww_x]_0^\pi+\frac A2[w^2]_0^\pi-K\int w_x^2+B\int w^2
=-K\int w_x^2+2BI\le0.\]
Vì $I(0)=0$, $w=0$. Phép đổi ẩn kiểm tra PDE, các sine kiểm tra biên, Fourier kiểm tra dữ liệu đầu trong miền. Dữ liệu không khớp Dirichlet tại các góc như đã thấy từ $f=1$ hoặc $f(\pi)=\pi$; nghiệm trơn khi $t>0$ và tiến đều về $0$.\end{document}